\chapter{The Semantic Embedding \texorpdfstring{$\iota : \mathcal{G} \hookrightarrow \mathcal{M}$}{iota : G ↪ M}}
\label{ch:semantic-embedding}

Part~I developed the intrinsic geometry of the RSVP manifold
$(\mathcal{M},g)$, including the mixed-type regimes, stratified
singularities, and harmonic attractors.  
Part~II developed the algebraic semantics of Polyxan hypergraphs:
typed atoms, curvature, sheaves, fibrations with connection, quine
structure, and cohomology.

The purpose of this chapter is to unify these constructions by defining
the \emph{semantic embedding}

\begin{equation}
\label{eq:iota-def}
\iota : \mathcal{G} \hookrightarrow \mathcal{M},
\end{equation}

a geometric realisation of the discrete semantic graph inside the
continuous semantic manifold.  
This embedding is not arbitrary: it is curvature-controlled,
entropy-aligned, and must satisfy a global energy minimisation
principle.  Its critical points are precisely the harmonic embeddings
that form the backbone of semantic galaxies.

\section{Incidence Complex and Stratification}
Let $\mathcal{G}$ be a finite typed Polyxan hypergraph.
Its geometric realisation is described by its \emph{incidence
complex}~$\mathrm{Inc}(\mathcal{G})$, a stratified simplicial set whose
$k$-simplices correspond to typed $k$-ary relations among content atoms.
Formally:

\begin{equation}
\mathrm{Inc}(\mathcal{G})
  = \coprod_{k \ge 0} \mathcal{G}_k \times \Delta^k,
\end{equation}

with face and degeneracy maps inherited from the hypergraph incidence
relations.  
This induces a stratification
$\bigsqcup S_k$ of the domain of~$\iota$.

The embedding map $\iota$ must respect this stratification:

\begin{equation}
\label{eq:stratified-embedding}
\iota(S_k) \subset \mathcal{M}^{(n-k)},
\end{equation}

where $\mathcal{M}^{(n-k)}$ is the $(n-k)$-stratum of the RSVP
stratification induced by curvature-critical surfaces (Chapter~13).
Thus higher-arity relations of $\mathcal{G}$ map to lower-dimensional
semantic strata.

\section{Type--Curvature Matching}

Let $\mathcal{H}_0[\mathcal{G}]$ be the primitive Polyxan curvature from
Chapter~21, defined as tag-entropy deficit plus hyperedge deficit.
Let $\mathcal{R}$ denote the scalar curvature of $(\mathcal{M},g)$.

The embedding~$\iota$ must satisfy the \emph{type--curvature matching
condition}:

\begin{equation}
\label{eq:type-curvature-matching}
\mathcal{H}_0[\mathcal{G}]
  = \int_{\iota(\mathrm{Inc}(\mathcal{G}))} \mathcal{R}\, dV_g.
\end{equation}

This balances the discrete curvature of $\mathcal{G}$ with the ambient
continuous curvature of $\mathcal{M}$.

It enforces that hypergraphs of higher curvature deficit occupy regions
of higher semantic curvature, and vice versa.

\section{Entropy Alignment}

Let $\mathcal{H}_{\mathrm{tag}}[\mathcal{G}]$ be the tag-sheaf entropy of
$\mathcal{G}$, defined as the Shannon entropy of tag distributions over
hyperedges.

The embedding must satisfy the \emph{entropy alignment constraint}:

\begin{equation}
\label{eq:entropy-alignment}
S\big\rvert_{\iota(\mathcal{G})}
  = \mathcal{H}_{\mathrm{tag}}[\mathcal{G}],
\end{equation}

i.e.\ the RSVP entropy field restricted to the image of~$\mathcal{G}$
matches the discrete tag entropy exactly.  
This guarantees semantic consistency between the continuous and discrete
representations of meaning.

\section{Embedding Energy Functional}

To couple both curvature and entropy, define the \emph{embedding
energy}:

\begin{equation}
\label{eq:embedding-energy}
\mathcal{E}[\iota]
  = \alpha \!\int_{\mathrm{Inc}(\mathcal{G})}
       \!\lVert \nabla \iota \rVert^2 \, d\mu
    + \beta \!\int_{\mathrm{Inc}(\mathcal{G})}
       \!\big(\mathcal{R} \circ \iota - \mathcal{H}_0[\mathcal{G}]\big)^2
       d\mu
    + \gamma \!\int_{\mathrm{Inc}(\mathcal{G})}
       \!\big(S \circ \iota - \mathcal{H}_{\mathrm{tag}}[\mathcal{G}]\big)^2
       d\mu.
\end{equation}

Here:
\begin{itemize}
\item The first term penalises geometric distortion.  
\item The second term enforces type--curvature matching.  
\item The third term enforces entropy alignment.
\end{itemize}

Critical points of~\eqref{eq:embedding-energy} satisfy the Euler--Lagrange
equations:

\begin{equation}
\label{eq:embedding-euler-lagrange}
\Delta \iota 
  = \beta\, \nabla\big(\mathcal{R}\circ\iota\big)
  + \gamma\, \nabla\big(S\circ\iota\big),
\end{equation}

which identifies harmonic embeddings as minimisers of~$\mathcal{E}[\iota]$.

\section{Harmonic Embeddings}

A semantic embedding is \emph{harmonic} if:

\begin{equation}
\label{eq:harmonic}
\Delta \iota = 0
\quad\text{and}\quad
\mathcal{R}\circ\iota = \mathcal{H}_0[\mathcal{G}],
\quad
S\circ\iota = \mathcal{H}_{\mathrm{tag}}[\mathcal{G}].
\end{equation}

Thus the image of $\mathcal{G}$ under $\iota$ lies entirely within an
RSVP stratum of matching curvature and entropy.

From Chapter~26, we know that Polyxan harmonic hypergraphs satisfy
$\Delta_{\mathcal{G}} \omega = 0$ for all cochains.  
Under~$\iota$, these discrete harmonic structures correspond exactly to
harmonic modes of the RSVP fields:

\begin{equation}
\iota^\ast(\text{RSVP harmonic}) 
  = 
\text{Polyxan harmonic}.
\end{equation}

Therefore harmonic embeddings identify the discrete and continuous
harmonic structures of meaning.

\section{Embedding Existence and Regularity Theorem}

We now prove the main result of this chapter.

\begin{theorem}[Embedding Existence and Regularity]
\label{thm:embedding-existence}
Every finite typed Polyxan hypergraph $\mathcal{G}$ admits a smooth,
stratified, type- and entropy-aligned embedding
$\iota : \mathcal{G} \hookrightarrow \mathcal{M}$ that minimises the
energy functional~\eqref{eq:embedding-energy}.  
In the elliptic region of $\mathcal{M}$, this embedding is unique up to
diffeomorphism.
\end{theorem}

\begin{proof}
Existence follows from the direct method of the calculus of variations:
$\mathcal{E}[\iota]$ is coercive, lower semicontinuous, and invariant
under reparameterisation of strata.  
A minimiser exists in $W^{1,2}$ and elliptic regularity implies smoothness.

Stratification compatibility follows from the coercivity of $\alpha$ and
the geometric constraints in~\eqref{eq:stratified-embedding}.

Type--curvature and entropy alignment follow from minimisation of the
$\beta$ and $\gamma$ terms, which force the matching conditions almost
everywhere, then everywhere by analyticity of $\mathcal{R}$ and $S$.

Uniqueness in the elliptic region follows from strict convexity of the
Dirichlet term.
\end{proof}

This establishes the fundamental bridge between Polyxan hypergraphs and
RSVP geometry.

\section{Flat Embedding of the Universal Media Quine}

Finally we identify the unique ground state of the coupled theory.

\begin{theorem}[Flat Embedding of $\mathfrak{U}$]
\label{thm:uinembedding}
The universal media quine $\mathfrak{U}$ embeds into the RSVP conformal
vacuum $(\Phi^\ast,0,S^\ast)$ via a unique (up to isometry) flat,
harmonic, stratified embedding $\iota^\ast$ whose image is a minimal
submanifold of $(\mathcal{M},g)$.  
This minimises the coupled action and realises the absolute ground state
of the RSVP--Polyxan theory.
\end{theorem}

\begin{proof}
Since $\mathfrak{U}$ has zero curvature deficit and vanishing positive-degree
cohomology (Chapter~26), the matching conditions enforce:

\[
\mathcal{R}\circ\iota^\ast = 0,
\qquad
S\circ\iota^\ast = S^\ast.
\]

Thus $\iota^\ast$ maps into the conformal RSVP vacuum.  
The harmonicity condition reduces to $\Delta \iota^\ast = 0$, implying
$\iota^\ast$ is a flat minimal immersion.  
Uniqueness follows from convexity of $\mathcal{E}$ under these
constraints and the rigidity of minimal embeddings into flat vacua.
\end{proof}

\medskip

Thus the universal quine $\mathfrak{U}$ and the RSVP conformal vacuum
are unified into a single, stable, harmonic ground state.

This completes the definition and analysis of the semantic embedding.
